% ═══════════════════════════════════════════════════════════════
%  chapters/ch05-adc-dma.tex
%  ADC + DMA: как захватить сигнал
% ═══════════════════════════════════════════════════════════════

\chapter{ADC и DMA: как захватить сигнал}
\label{ch:adc-dma}

\epigraph{<<ADC~--- это переводчик с аналогового языка на цифровой.
DMA~--- это курьер, который носит данные без участия CPU>>.}{}

\section{Введение}
\label{sec:adc-intro}

В \texttt{ch04} мы разобрали \emph{протокол}~--- язык, на
котором \emph{общаются} ПК и платы. Но \emph{что} именно
передаёт осциллограф? \textbf{Сигнал}. А сигнал~---
это \emph{напряжение}, которое \emph{меняется} во времени.

\textbf{Задача} осциллографа~--- \emph{захватить} этот
сигнал \emph{как можно точнее} и \emph{как можно быстрее}.
Для этого нужны \emph{два} инструмента: \textbf{ADC}
и \textbf{DMA}.

\begin{ideabox}
\textbf{Почему это \emph{сложно}?} Представьте, что вы
пытаетесь \emph{сфотографировать} \emph{летящую} стрелу.
Обычная камера \emph{размажет} её~--- слишком медленная
выдержка. Нужна \emph{скоростная} камера~--- чтобы
\emph{застыло} каждое мгновение. ADC~--- это \emph{скоростная}
камера для \emph{электричества}.
\end{ideabox}

\subsection{Что вы узнаете}

\begin{itemize}
    \item Что такое ADC и зачем он нужен.
    \item Как работает \emph{семплирование} (дискретизация).
    \item Какие параметры важны: разрядность, частота, опора.
    \item Что такое DMA и \emph{зачем} он нужен.
    \item Как работает \emph{двойной буфер} (ping-pong).
    \item Почему \emph{прерывания}~--- это ключ к \emph{непрерывному}
          захвату.
\end{itemize}

\section{Что такое ADC}
\label{sec:adc-what}

\textbf{ADC} (Analog-to-Digital Converter, аналого-цифровой
преобразователь)~--- это \emph{микросхема}, которая
\emph{превращает} \emph{аналоговое} напряжение в
\emph{цифровое} число.

\begin{keybox}
\textbf{Аналогия: переводчик.} Аналоговый сигнал~--- это
\emph{плавная} кривая (как \emph{голос} человека).
Цифровой~--- это \emph{набор} чисел (как \emph{текст}).
ADC~--- это \emph{переводчик}, который \emph{слушает}
голос и \emph{записывает} его \emph{словами}.
\end{keybox}

\subsection{Как работает ADC}

\textbf{Представьте}, что у вас есть \emph{линейка} от
\num{0} до \num{3.3} вольт. ADC \emph{измеряет} напряжение
и \emph{округляет} его до \emph{ближайшего} деления.

Если у ADC \textbf{12 бит}, то у него \textbf{4096}
делений (от \num{0} до \num{4095}). Значит, каждое деление~---
это:
\[
\frac{3.3 \text{ В}}{4096} \approx 0.8 \text{ мВ}
\]
То есть ADC \emph{различает} напряжение с точностью до
\textbf{0.8 милливольта}.

\begin{center}
\begin{tabular}{|c|c|}
\hline
\textbf{Напряжение} & \textbf{Число (12 бит)} \\
\hline
\num{0.0} В & \num{0} \\
\num{1.65} В & \num{2048} \\
\num{3.3} В & \num{4095} \\
\hline
\end{tabular}
\end{center}

\subsection{Семплирование (дискретизация)}

ADC не \emph{непрерывен}. Он \emph{измеряет} напряжение
\emph{через} \emph{равные} промежутки времени. Этот процесс
называется \textbf{семплированием} (или \emph{дискретизацией}).

\textbf{Пример}: если частота дискретизации~--- \SI{100}{\kilo\hertz}
(\num{100000} семплов в секунду), то ADC \emph{измеряет}
напряжение \emph{каждые} \SI{10}{\micro\second}.

\begin{ideabox}
\textbf{Аналогия: киносъёмка.} Фильм~--- это \emph{24}
кадра в секунду. Наш глаз \emph{не} видит \emph{между}
кадрами~--- мозг \emph{склеивает} их в \emph{плавное}
движение. ADC работает так же: \emph{много} семплов
в секунду~--- и \emph{кажется}, что сигнал \emph{непрерывен}.
\end{ideabox}

\subsection{Частота дискретизации и теорема Котельникова}

\textbf{Ключевой вопрос}: \emph{как часто} надо семплировать,
чтобы \emph{не потерять} сигнал?

\textbf{Теорема Котельникова} (или Найквиста): частота
дискретизации должна быть \emph{как минимум в два раза}
выше \emph{максимальной} частоты сигнала.

\begin{keybox}
\textbf{Пример.} Если сигнал~--- \SI{1}{\kilo\hertz} (синус),
то семплировать надо \emph{минимум} \SI{2}{\kilo\hertz}.
Но \emph{на практике}~--- в \emph{10} раз чаще: \SI{10}{\kilo\hertz}.
Иначе \emph{форма} сигнала будет \emph{искажена}.
\end{keybox}

\begin{warnbox}
\textbf{Что такое <<алиасинг>>?} Если семплировать \emph{слишком}
редко~--- быстрый сигнал \emph{превратится} в \emph{медленный}.
Например, синус \SI{1}{\kilo\hertz} при семплировании
\SI{1.5}{\kilo\hertz} \emph{выглядит} как \SI{500}{\hertz}.
Это~--- \emph{обман}. Читатель \emph{думает}, что видит
один сигнал, а на самом деле~--- \emph{другой}.
\end{warnbox}

\section{ADC на STM32F401}
\label{sec:adc-stm32}

Наш осциллограф~--- на \textbf{NUCLEO-F401RE}. Внутри~---
\textbf{ADC} со следующими параметрами:

\begin{itemize}
    \item \textbf{Разрядность}: \SI{12}{\bit} (4096 уровней).
    \item \textbf{Максимальная частота}: \SI{2.4}{\mega\sample/\second}
      (\textbf{2.4 MSa/s}~--- мегасемплов в секунду).
    \item \textbf{Опорное напряжение}: \SI{3.3}{\volt} (по умолчанию).
    \item \textbf{Каналы}: \num{16} (разные пины).
\end{itemize}

\begin{notebox}
\textbf{Что такое \texttt{Sa/s}?} \texttt{Sa/s} (samples per second)~---
это \emph{семплов в секунду}. Единица измерения \emph{частоты
дискретизации}. \textbf{Почему не Гц?} Потому что \emph{герц}~---
это <<циклов в секунду>>. ADC \emph{не} делает циклы~---
он делает \emph{измерения}. Чтобы не путать~--- используют
\texttt{Sa/s}. Например, \texttt{100 kSa/s}~--- это
\emph{100000} семплов в секунду.
\end{notebox}

\begin{keybox}
\textbf{Почему \SI{12}{\bit}, а не \SI{16}{\bit}?} \SI{12}{\bit}~---
это \emph{стандарт} для STM32. \SI{16}{\bit}~--- \emph{дороже}
и \emph{медленнее}. Для \emph{учебного} осциллографа
\SI{12}{\bit}~--- \emph{достаточно}.
\end{keybox}

\subsection{Пины и каналы}

На Nucleo-F401RE \emph{аналоговые} входы~--- это пины
\texttt{PA0}--\texttt{PA7}, \texttt{PB0}--\texttt{PB1},
\texttt{PC0}--\texttt{PC5}.

\textbf{Мы используем} \texttt{PA0}~--- это \textbf{ADC1\_IN0}.
К нему \emph{подключим} генератор (или другой источник
сигнала).

\subsection{Опорное напряжение}

ADC \emph{сравнивает} входное напряжение с \textbf{опорным}
(\texttt{VREF+}). На Nucleo-F401RE \texttt{VREF+}~---
это \SI{3.3}{\volt}.

\textbf{Что это значит}: если подать \SI{3.3}{\volt}~---
ADC вернёт \num{4095}. Если \SI{0}{\volt}~--- \num{0}.
Если \SI{1.65}{\volt}~--- \num{2048}.

\begin{warnbox}
\textbf{Не подавайте больше \SI{3.3}{\volt}!} ADC \emph{сгорит}
(или \emph{повредится}). Если сигнал \emph{больше}~---
нужен \emph{делитель} напряжения.
\end{warnbox}

\section{Зачем нужен DMA}
\label{sec:dma-why}

\textbf{Представьте}: ADC \emph{семплирует} сигнал с
частотой \SI{100}{\kilo\hertz}. Каждый семпл~--- это
\emph{число} (\SI{12}{\bit}). \emph{Кто-то} должен
\emph{забрать} это число \emph{из} регистра ADC и
\emph{положить} в \emph{память} (SRAM).

\subsection{Плохой вариант: CPU}

\textbf{CPU} мог бы \emph{читать} регистр ADC \emph{каждый}
раз, когда готов \emph{новый} семпл. Но:

\begin{itemize}
    \item \textbf{CPU занят}~--- не может делать \emph{ничего}
          другого.
    \item \textbf{Прерывания}~--- \num{100000} прерываний
          в секунду. \emph{Много}.
    \item \textbf{Потери}~--- если CPU \emph{опоздал},
          семпл \emph{потерян}.
\end{itemize}

\begin{warnbox}
\textbf{Аналогия: курьер и директор.} Директор (CPU)
\emph{сам} бегает \emph{за} каждой \emph{посылкой}~---
это \emph{глупо}. Директор \emph{должен} \emph{думать},
а \emph{не} \emph{носить} коробки.
\end{warnbox}

\subsection{Хороший вариант: DMA}

\textbf{DMA} (Direct Memory Access)~--- это \emph{отдельный}
контроллер, который \emph{умеет} \emph{перекладывать}
данные \emph{из} периферии \emph{в} память \emph{без}
участия CPU.

\textbf{Как это работает}:
\begin{enumerate}
    \item CPU \emph{настраивает} DMA: <<когда ADC
          \emph{готов}~--- \emph{забери} число и \emph{положи}
          в \emph{этот} адрес>>.
    \item CPU \emph{забывает} про ADC. \emph{Занимается}
          другими делами.
    \item ADC \emph{семплирует}. Когда готов~--- \emph{сообщает}
          DMA.
    \item DMA \emph{забирает} число и \emph{кладёт} в память.
    \item CPU \emph{даже не знает}, что это произошло.
\end{enumerate}

\begin{keybox}
\textbf{DMA~--- это <<курьер>>.} Он \emph{носит} данные
\emph{без} участия CPU. CPU \emph{только} \emph{настраивает}
его~--- \emph{один} раз. Дальше~--- \emph{всё} \emph{само}.
\end{keybox}

\subsection{Почему DMA критичен для осциллографа}

\textbf{Без DMA} CPU \emph{не успевает} обрабатывать
\SI{100}{\kilo\sample/\second}. Он \emph{тратит} всё
время на \emph{чтение} ADC.

\textbf{С DMA} CPU \emph{свободен}~--- он \emph{занимается}
USB, \emph{отправкой} данных, \emph{обработкой} команд.
\emph{Только} DMA \emph{работает} с ADC.

\section{Двойной буфер (ping-pong)}
\label{sec:dma-pingpong}

\textbf{Проблема}: DMA \emph{пишет} в память. CPU \emph{читает}
из памяти. Что если CPU \emph{читает} \emph{ту} ячейку,
в \emph{которую} DMA \emph{сейчас} \emph{пишет}?

\textbf{Решение}: \textbf{двойной буфер} (ping-pong). Память
\emph{делится} на \emph{две} половины. DMA \emph{пишет}
в \emph{одну}, CPU \emph{читает} из \emph{другой}. \emph{Потом}
они \emph{меняются} местами.




\begin{figure}[htbp]
\centering
\adjustbox{max width=\textwidth}{%
\begin{tikzpicture}[
    x=1cm, y=1cm,
    buffer/.style={
        draw, thick,
        minimum width=4cm,
        minimum height=1cm,
        align=center,
        font=\small\ttfamily,
        fill=white
    },
    active/.style={buffer, fill=clrHWData!30},
    free/.style={buffer, fill=clrLogicData!20},
    label/.style={
        font=\footnotesize,
        above,
        inner sep=2pt
    },
    phase/.style={
        font=\small\bfseries,
        fill=white,
        inner sep=3pt
    }
]

% ─── Координатная сетка ────────────────────────────
\draw[step=1cm, gray!20, very thin] (-3, -3) grid (10, 3);

% ─── Фаза 1 ───────────────────────────────────────
\node[phase] at (0, 2) {Фаза 1};
\node[active] (buf1) at (1, 1) {Буфер 1 [DMA пишет]};
\node[free]   (buf2) at (6, 1) {Буфер 2 [CPU читает]};

% Стрелки
\draw[-{Latex[length=3mm]}, very thick, color=clrFlowData]
    (-2, 1) -- (buf1.west);
\node[font=\footnotesize, left] at (-2.1, 1) {ADC};

\draw[-{Latex[length=3mm]}, very thick, color=clrFlowCmd]
    (buf2.east) -- (9, 1);
\node[font=\footnotesize, right] at (9.0, 1) {USB};

% ─── Разделитель ──────────────────────────────────
\draw[dashed, gray!90] (-3, 0) -- (10, 0);

% ─── Фаза 2 ───────────────────────────────────────
\node[phase] at (0, -0.5) {Фаза 2};
\node[free]   (buf3) at (1, -1.5) {Буфер 1 [CPU читает]};
\node[active] (buf4) at (7.5, -1.5) {Буфер 2 [DMA пишет]};

% Стрелки
\draw[-{Latex[length=3mm]}, very thick, color=clrFlowData]
    (4.5, -1.5) -- (buf4.west);
\node[font=\footnotesize, above] at (4, -1.7) {ADC};

\draw[-{Latex[length=3mm]}, very thick, color=clrFlowCmd]
    (buf3.west) -- (-2, -1.5);
\node[font=\footnotesize, left] at (-2.2, -1.5) {USB};

% ─── Метки осей (для отладки) ─────────────────────
\foreach \x in {-3,-2,-1,0,1,2,3,4,5,6,7,8,9,10} {
    \node[font=\tiny, gray] at (\x, -3.3) {\x};
}
\foreach \y in {-3,-2,-1,0,1,2,3} {
    \node[font=\tiny, gray] at (-3.3, \y) {\y};
}

\end{tikzpicture}
}% end adjustbox
\caption{Двойной буфер (ping-pong). В фазе 1 DMA пишет
в буфер 1, CPU читает из буфера 2. В фазе 2 они меняются
местами.}
\label{fig:dma-pingpong}
\end{figure}



\textbf{Как это работает}:
\begin{enumerate}
    \item \textbf{Фаза 1}: DMA \emph{пишет} в \emph{первую}
          половину. CPU \emph{читает} из \emph{второй}.
    \item \textbf{Фаза 1 закончена}: DMA \emph{заполнил}
          первую половину. Срабатывает \emph{прерывание}
          (\texttt{half-cplt}).
    \item \textbf{Фаза 2}: DMA \emph{пишет} во \emph{вторую}
          половину. CPU \emph{читает} из \emph{первой}.
    \item \textbf{Фаза 2 закончена}: DMA \emph{заполнил}
          вторую половину. Срабатывает \emph{прерывание}
          (\texttt{cplt}).
    \item \textbf{Повтор} с шага 1.
\end{enumerate}

\begin{keybox}
\textbf{Почему это работает?} DMA \emph{никогда} не
пишет в ту половину, которую CPU \emph{сейчас} читает.
Они \emph{всегда} в \emph{разных} половинах. \emph{Потери}
данных~--- \emph{нет}.
\end{keybox}

\begin{warnbox}
\textbf{Что если CPU \emph{опоздал}?} DMA \emph{уже}
заполнил \emph{свою} половину и \emph{переключился}
на \emph{ту}, которую CPU \emph{читает}. \emph{Данные}
\emph{потеряны}. \emph{Поэтому} CPU \emph{обязан}
обработать половину \emph{быстро}~--- за \emph{время}
заполнения \emph{одной} половины.
\end{warnbox}




%%%%%%=====================



\section{Прерывания: сигнал <<готово>>}
\label{sec:adc-interrupts}

В \texttt{sec:dma-pingpong} мы увидели \emph{двойной}
буфер. Но \emph{как} CPU \emph{узнаёт}, что \emph{половина}
заполнена? Через \textbf{прерывания}.

\subsection{Что такое прерывание}

\textbf{Прерывание} (interrupt)~--- это \emph{сигнал}
от \emph{периферии} к \emph{процессору}: <<\emph{остановись},
у меня \emph{важное} дело>>.

\begin{ideabox}
\textbf{Аналогия: звонок в дверь.} Вы \emph{работаете}
(CPU). Вдруг~--- \emph{звонок} (прерывание). Вы
\emph{открываете} дверь, \emph{принимаете} посылку
(обработчик), \emph{возвращаетесь} к работе. \emph{Работа}
\emph{не} потеряна~--- она \emph{продолжается} с того
же места.
\end{ideabox}

\textbf{Как это работает}:
\begin{enumerate}
    \item \textbf{Периферия} (DMA) \emph{выставляет} флаг
          <<я готов>>.
    \item \textbf{Контроллер} \emph{прерываний} (NVIC)
          \emph{видит} флаг и \emph{сообщает} CPU.
    \item \textbf{CPU} \emph{сохраняет} текущее состояние
          (регистры) в \emph{стек}.
    \item \textbf{CPU} \emph{переходит} в \emph{обработчик}
          прерывания (ISR~--- Interrupt Service Routine).
    \item \textbf{Обработчик} \emph{делает} своё дело
          и \emph{возвращает} управление.
    \item \textbf{CPU} \emph{восстанавливает} состояние
          из стека и \emph{продолжает} работу.
\end{enumerate}

\begin{keybox}
\textbf{Прерывание~--- это \emph{асинхронность}.} CPU
\emph{не} ждёт~--- он \emph{занимается} своими делами.
Когда \emph{периферия} готова~--- она \emph{сообщает}.
CPU \emph{реагирует} \emph{мгновенно}.
\end{keybox}

\subsection{Два прерывания DMA: half-cplt и cplt}

\textbf{DMA} \emph{генерирует} \emph{два} прерывания на
\emph{цикл} двойного буфера:

\begin{itemize}
    \item \texttt{HAL\_ADC\_ConvHalfCpltCallback()}~---
          \emph{первая} половина заполнена.
    \item \texttt{HAL\_ADC\_ConvCpltCallback()}~---
          \emph{вторая} половина заполнена.
\end{itemize}

\textbf{Схема}:
\begin{enumerate}
    \item DMA \emph{пишет} в буфер 1. CPU \emph{читает}
          буфер 2.
    \item DMA \emph{заполнил} буфер 1~--- \emph{прерывание}
          \texttt{half-cplt}. CPU \emph{переключается}:
          \emph{начинает} читать буфер 1.
    \item DMA \emph{пишет} в буфер 2. CPU \emph{читает}
          буфер 1.
    \item DMA \emph{заполнил} буфер 2~--- \emph{прерывание}
          \texttt{cplt}. CPU \emph{переключается}:
          \emph{начинает} читать буфер 2.
    \item Повтор с шага 1.
\end{enumerate}

\begin{warnbox}
\textbf{Обработчик прерывания должен быть \emph{быстрым}.}
Пока CPU \emph{в} обработчике~--- он \emph{не} делает
\emph{других} дел. Если обработчик \emph{медленный}~---
\emph{пропускаются} другие прерывания. \textbf{Правило}:
обработчик~--- \emph{минимум} работы. \emph{Тяжёлое}~---
в \emph{основной} цикл.
\end{warnbox}

\section{Кольцо кадров}
\label{sec:adc-ring}

\textbf{Прерывание} даёт \emph{сигнал}: <<половина готова>>.
Но \emph{что} делать с \emph{данными}? Их надо \emph{собрать}
в \emph{кадр} для \emph{отправки} на ПК.

\textbf{Кольцо кадров}~--- это \emph{буфер} \emph{кадров}.
Пока DMA \emph{пишет} в \emph{один} кадр~--- CPU \emph{отправляет}
\emph{другой}.

\begin{figure}[htbp]
\centering
\adjustbox{max width=\textwidth}{%
\begin{tikzpicture}[
    x=1cm, y=1cm,
    box/.style={
        draw, thick, rounded corners=1mm,
        minimum width=2cm, minimum height=1cm,
        align=center, font=\footnotesize\ttfamily,
        fill=white
    },
    wr/.style={box, fill=clrHWData!30},
    rd/.style={box, fill=clrUIData!20},
    fre/.style={box, fill=gray!10},
    arr/.style={-{Latex[length=3mm]}, very thick}
]

% Кольцо из 4 кадров
\node[wr]  (f0) at (0, 0) {Кадр 0\\[2pt]\textbf{DMA}};
\node[fre] (f1) at (3, 0) {Кадр 1\\[2pt]пусто};
\node[fre] (f2) at (6, 0) {Кадр 2\\[2pt]пусто};
\node[rd]  (f3) at (9, 0) {Кадр 3\\[2pt]\textbf{USB}};

% Стрелки (по кругу)
\draw[arr] (f0.east) -- (f1.west);
\draw[arr] (f1.east) -- (f2.west);
\draw[arr] (f2.east) -- (f3.west);
\draw[arr, dashed] (f3.south) to[bend left=40] (f0.south);

% Подписи
\node[font=\footnotesize, above=0.3cm of f0] {пишется};
\node[font=\footnotesize, above=0.3cm of f3] {отправляется};

\end{tikzpicture}
}% end adjustbox
\caption{Кольцо кадров. DMA пишет в текущий кадр,
USB отправляет уже готовый. Когда DMA заканчивает~---
переходит к следующему.}
\label{fig:adc-ring}
\end{figure}

\textbf{Как это работает}:
\begin{enumerate}
    \item DMA \emph{заполняет} кадр 0.
    \item Когда кадр 0 \emph{готов}~--- DMA \emph{переходит}
          к кадру 1. \emph{Кадр 0} становится \emph{доступен}
          для \emph{отправки}.
    \item USB \emph{отправляет} кадр 0.
    \item DMA \emph{заполняет} кадр 1. USB \emph{отправляет}
          кадр 1 (когда \emph{готов}).
    \item \emph{Повтор} по кругу.
\end{enumerate}

\begin{keybox}
\textbf{Почему кольцо?} Потому что \emph{непрерывно}.
DMA \emph{никогда} не \emph{останавливается}. Он
\emph{пишет} в \emph{следующий} кадр. Пока USB
\emph{отправляет} \emph{предыдущий}.
\end{keybox}

\begin{warnbox}
\textbf{Что если USB \emph{опоздал}?} DMA \emph{обогнал}
USB \emph{на круг}~--- и \emph{пишет} в тот кадр,
который USB \emph{ещё} \emph{отправляет}. \emph{Данные}
\emph{испорчены}. \textbf{Решение}: \emph{увеличить}
размер кольца или \emph{ускорить} USB.
\end{warnbox}

\section{Код: adc.c и dma.c}
\label{sec:adc-code}

\textbf{Теперь}~--- \emph{код}. Он \emph{короткий}, но
\emph{плотный}. \emph{Каждая} строка \emph{важна}.

\subsection{Настройка ADC}

\begin{listing}[htbp]
\begin{minted}[frame=none, bgcolor=gray!5]{c}
void adc_init(void)
{
    // 1. Тактирование ADC1
    RCC->APB2ENR |= RCC_APB2ENR_ADC1EN;

    // 2. Настройка ADC1
    ADC1->CR1 = 0;                  // 12-bit, no scan
    ADC1->CR2 = ADC_CR2_ADON;       // включить ADC
    ADC1->SQR1 = 0;                 // 1 преобразование
    ADC1->SQR3 = 0;                 // канал 0 (PA0)

    // 3. Время семплирования
    ADC1->SMPR2 = (7 << 0);         // 480 циклов

    // 4. Калибровка
    ADC1->CR2 |= ADC_CR2_RSTCAL;
    while (ADC1->CR2 & ADC_CR2_RSTCAL);
    ADC1->CR2 |= ADC_CR2_CAL;
    while (ADC1->CR2 & ADC_CR2_CAL);

    // 5. Включить ADC
    ADC1->CR2 |= ADC_CR2_ADON;
}
\end{minted}
\caption{\texttt{adc\_init()}~--- настройка ADC1}
\label{lst:adc-init}
\end{listing}

\textbf{Что здесь важно}:
\begin{itemize}
    \item \texttt{RCC->APB2ENR}~--- \emph{включает} тактирование
          ADC. \emph{Без этого} ADC \emph{не работает}.
    \item \texttt{ADC1->CR1 = 0}~--- \SI{12}{\bit}, без
          \emph{сканирования} (один канал).
    \item \texttt{ADC1->SQR3 = 0}~--- канал \texttt{0}
          (\texttt{PA0}).
    \item \texttt{ADC1->SMPR2 = (7 << 0)}~--- \emph{время}
          семплирования. \texttt{7}~--- \emph{максимум} (\num{480}
          циклов). \emph{Точнее}, но \emph{медленнее}.
    \item \textbf{Калибровка}~--- \emph{обязательна} перед
          \emph{первым} запуском.
\end{itemize}

\subsection{Настройка DMA}

\begin{listing}[htbp]
\begin{minted}[frame=none, bgcolor=gray!5]{c}
void dma_init(uint16_t *buf, size_t len)
{
    // 1. Тактирование DMA2
    RCC->AHB1ENR |= RCC_AHB1ENR_DMA2EN;

    // 2. Поток 0, канал 0 (ADC1)
    DMA2_Stream0->CR = 0;
    DMA2_Stream0->PAR = (uint32_t)&ADC1->DR;
    DMA2_Stream0->M0AR = (uint32_t)buf;
    DMA2_Stream0->NDTR = len;

    // 3. Настройки
    DMA2_Stream0->CR =
        (0 << 25)  |   // канал 0
        (1 << 13)  |   // memory increment
        (1 << 11)  |   // 16-bit peripheral
        (1 << 10)  |   // 16-bit memory
        (1 << 8)   |   // circular mode
        (1 << 4);      // half-transfer interrupt

    // 4. Включить
    DMA2_Stream0->CR |= DMA_SxCR_EN;
    ADC1->CR2 |= ADC_CR2_DMA;
}
\end{minted}
\caption{\texttt{dma\_init()}~--- настройка DMA}
\label{lst:dma-init}
\end{listing}

\textbf{Что здесь важно}:
\begin{itemize}
    \item \texttt{PAR}~--- \emph{адрес} периферии
          (\texttt{ADC1->DR}~--- \emph{регистр} данных).
    \item \texttt{M0AR}~--- \emph{адрес} памяти (буфер).
    \item \texttt{NDTR}~--- \emph{сколько} передать.
    \item \texttt{(1 << 8)}~--- \emph{циклический} режим.
          Когда \emph{всё} передано~--- \emph{начинает}
          заново.
    \item \texttt{(1 << 4)}~--- \emph{прерывание} на
          \emph{половине} (\texttt{half-cplt}).
\end{itemize}

\subsection{Обработчики прерываний}

\begin{listing}[htbp]
\begin{minted}[frame=none, bgcolor=gray!5]{c}
void HAL_ADC_ConvHalfCpltCallback(ADC_HandleTypeDef *hadc)
{
    // Первая половина готова
    push_block(&dma_buf[0], OSC_DMA_POINTS);
}

void HAL_ADC_ConvCpltCallback(ADC_HandleTypeDef *hadc)
{
    // Вторая половина готова
    push_block(&dma_buf[OSC_DMA_POINTS], OSC_DMA_POINTS);
}
\end{minted}
\caption{Обработчики прерываний DMA}
\label{lst:dma-irq}
\end{listing}

\textbf{Что здесь важно}:
\begin{itemize}
    \item \texttt{HAL\_ADC\_ConvHalfCpltCallback}~--- \emph{первая}
          половина \texttt{dma\_buf} готова.
    \item \texttt{HAL\_ADC\_ConvCpltCallback}~--- \emph{вторая}
          половина.
    \item \texttt{push\_block()}~--- \emph{копирует}
          данные в \emph{кольцо} кадров.
\end{itemize}

\begin{warnbox}
\textbf{Обработчик должен быть \emph{коротким}.} В
\emph{идеале}~--- только \texttt{memcpy()}. \emph{Всё}
остальное~--- в \emph{основной} цикл. \emph{Почему}?
Пока CPU \emph{в} обработчике~--- он \emph{не} делает
\emph{других} дел. \emph{Долгий} обработчик~---
\emph{пропуск} других прерываний.
\end{warnbox}

\section{Уроки главы}
\label{sec:adc-lessons}

\begin{enumerate}
    \item \textbf{ADC~--- это переводчик.} Он \emph{превращает}
          напряжение в \emph{число}.
    \item \textbf{12 бит~--- 4096 уровней.} Точность~---
          \SI{0.8}{\milli\volt}.
    \item \textbf{Семплирование~--- это <<фотографирование>>}
          сигнала \emph{через} \emph{равные} промежутки.
    \item \textbf{Теорема Котельникова:} частота семплирования
          \emph{минимум} \emph{вдвое} выше \emph{максимальной}
          частоты сигнала.
    \item \textbf{DMA~--- это курьер.} \emph{Носит} данные
          \emph{без} участия CPU.
    \item \textbf{Ping-pong~--- это двойной буфер.} Пока DMA
          \emph{пишет} в \emph{одну} половину~--- CPU \emph{читает} из \emph{другой}. Ping-pong \emph{не} теряет данные.
    \item \textbf{Прерывания~--- \emph{сигнал}, что половина
      \emph{готова}.} \emph{Быстрые}
          обработчики~--- \emph{залог} скорости.
    \item \textbf{Кольцо кадров~--- непрерывность.} DMA
          \emph{пишет}, USB \emph{отправляет}.
    \item \textbf{Код~--- плотный.} \emph{Каждая} строка
          \emph{важна}.
\end{enumerate}


\section{Чек-лист грабель}
\label{sec:adc-pitfalls}

\begin{table}[htbp]
\centering
\caption{Грабли ADC и DMA}
\label{tab:adc-pitfalls}
\small
\begin{tabularx}{\textwidth}{@{}c L L@{}}
\toprule
\textbf{№} & \textbf{Грабля} & \textbf{Решение} \\
\midrule
1 & ADC не работает & Включить \texttt{RCC\_APB2ENR\_ADC1EN} \\
2 & DMA не работает & Включить \texttt{RCC\_AHB1ENR\_DMA2EN} \\
3 & Неправильный канал & \texttt{ADC1->SQR3 = 0} для \texttt{PA0} \\
4 & Нет калибровки & \texttt{ADC\_CR2\_CAL} \\
5 & Медленное семплирование & \texttt{SMPR2 = 7} (\num{480} циклов) \\
6 & Потеря данных & Ping-pong + кольцо кадров \\
7 & Долгий обработчик & Только \texttt{memcpy()}, \emph{больше} ничего \\
8 & Ошибка в \texttt{NDTR} & \texttt{len} = размер \emph{одной} половины \\
\bottomrule
\end{tabularx}
\end{table}

\section{Упражнения}
\label{sec:adc-exercises}

\begin{exercisebox}
\textbf{Упражнение 1.} Что произойдёт, если \emph{не}
включить \texttt{RCC\_APB2ENR\_ADC1EN}? Почему?
\end{exercisebox}

\begin{exercisebox}
\textbf{Упражнение 2.} Нарисуйте \emph{диаграмму} кольца
кадров для \emph{четырёх} кадров. Что произойдёт, если
USB \emph{отстаёт} на \emph{два} кадра?
\end{exercisebox}

\needspace{5\baselineskip}
\begin{exercisebox}
\textbf{Упражнение 3.} Сколько времени занимает \emph{одно}
преобразование ADC при \texttt{SMPR2 = 7} (\num{480} циклов)
и \emph{частоте} ADC \SI{84}{\mega\hertz}? (\emph{Подсказка}:
время = \texttt{480 / 84e6}.)
\end{exercisebox}

\begin{exercisebox}
\textbf{Упражнение 4.} Предложите способ \emph{ускорить}
обработчик \texttt{push\_block()}, если данных \emph{много}.
Как \emph{избежать} \texttt{memcpy()}?
\end{exercisebox}

\section{Проверьте сами}
\label{sec:adc-check}

\begin{enumerate}
    \item Подключите Nucleo к компьютеру.
    \item Соберите прошивку (см. главу 2).
    \item Подключите генератор к \texttt{PA0} (\SI{1}{\kilo\hertz},
          \SI{1}{\volt}pp).
    \item Запустите OpenOCD + GDB.
    \item Поставьте точку останова в \texttt{HAL\_ADC\_ConvHalfCpltCallback}.
    \item Посмотрите, \emph{как} заполняется буфер.
    \item Сделайте \emph{скриншот} — \emph{пришлите} в книгу.
\end{enumerate}

\begin{ideabox}
\textbf{Если получилось~--- вы поняли ADC и DMA.} Это
\emph{главное}. Без них \emph{нет} осциллографа. \emph{Есть}
только \emph{железо}. Дальше~--- \emph{триггер},
\emph{кольцо}, \emph{USB}. \emph{Всё} \emph{впереди}.
\end{ideabox}



